\section{Equivalent layer classes}
\label{sec:eqlayer_classes}

\subsection{Dirichlet's equivalent layers}

\begin{definition}
	An equivalent layer $p_{\varepsilon}$ (Definition \ref{def:continuous-eqlayer}) is called Dirichlet's equivalent 
	layer if it has the following form in the Fourier domain
	\begin{equation}
		\mathcal{F} [ p_{\varepsilon}(\mathbf{r}'') ] =
		\alpha_{d} \, \mathcal{F} [ U_{\varepsilon}(\mathbf{r}'') ]  \: .
		\label{eq:p-Fourier-Dirichlet}
	\end{equation}
	\label{def:eql-Dirichlet}
\end{definition}
Note that Dirichlet's layers (Definition \ref{def:eql-Dirichlet}) are obtained from 
equation \ref{eq:p-Fourier} by setting $\alpha_{a} \equiv \alpha(1, \hat{\mathbf{z}}) = k$ 
(equation \ref{eq:diff-operators-generic}).
In this case, we can easily verify from equations \ref{eq:observed-data-Fourier},
\ref{eq:potential-U-epsilon-Fourier-layer-1} and \ref{eq:potential-U-epsilon-Fourier-layer-2} that they 
approximate the downward continuation of the observed potential field $d^{o}$ to the layer plane $\mathcal{S}''$.
It is important noting that such layers depend on $\alpha_{d}$, which 
carries information about the observed potential field (equation \ref{eq:observed-data-Fourier}).
In this case, the corresponding predicted field (equation \ref{eq:predicted-data-continuous-layer}) assumes the form
\begin{equation}
	d_{\varepsilon}(\mathbf{r})
	= \int\limits_{-\infty}^{\infty}\int\limits_{-\infty}^{\infty}
	p_{\varepsilon}(\mathbf{r}'') \: \tfrac{\partial}{\partial z} \Psi(\mathbf{r}, \mathbf{r}'') \:\: d\mathcal{S}''
	\: , \quad z < z''_{0} \: .
	\label{eq:predicted-data-Dirichlet-layer}
\end{equation}
From the physical point of view, this $d_{\varepsilon}$ is analogous
to the vertical component of the gravitational attraction due to a layer of monopoles
or the magnetic scalar potential due to a layer of vertically-magnetized dipoles.
Moreover, equation \ref{eq:predicted-data-Dirichlet-layer} approximates the well-known \textit{upward continuation integral} 
\citep[e.g.,][]{skeels1947, peters1949, henderson-zietz1949, dean1958, roy1962, bhattacharyya1967, henderson1970},
which is traditionally derived from the solution of the Dirichlet’s problem on the
layer plane (equation \ref{eq:U-Dirichlet-Neumann} with $\ell = 0$). 

Dirichlet's equivalent layers are the basis of the pioneer method proposed by \cite{dampney1969} for processing gravity data.
Such layers lead to potential-field transformations (equations 
\ref{eq:transformed-field-continuous-layer} and \ref{eq:transformed-field-continuous-layer-Fourier}) defined by 
$\alpha_{t} = \nicefrac{k \, \alpha_{f}}{\alpha_{d}}$ (equation \ref{eq:alpha-t-generic}).
Dirichlet's equivalent layers are commonly used to perform upward/downward continuation of the observed potential 
field, which is obtained by setting $\alpha_{f} = \alpha_{d}$, resulting in $\alpha_{t} = k$. 
Other potential-field transformations can also be accomplished from such equivalent layers, but they must be 
defined in terms of $\alpha_{f} \ne \alpha_{d}$.
In this case, however, defining the resulting transformed field $f_{\varepsilon}$ explicitly in the space domain 
(equation \ref{eq:transformed-field-continuous-layer}) 
from its counterpart in the Fourier domain (equation \ref{eq:transformed-field-continuous-layer-Fourier}) may be very difficult.


\subsection{Equivalent layers having minimum phase changes}

\begin{definition}
	An equivalent layer $p_{\varepsilon}$ (Definition \ref{def:continuous-eqlayer}) is said to have minimum 
	phase changes (with respect to the spatial distribution of the true sources) if its phase spectrum is given by
	\begin{equation}
		\Theta [ p_{\varepsilon}(\mathbf{r}'') ] = \Theta [ U_{\varepsilon}(\mathbf{r}^{\dagger}) ] \: ,
		\label{eq:p-phase-spectrum-U}
	\end{equation}
	where $\Theta [ U_{\varepsilon}(\mathbf{r}^{\dagger}) ]$ denotes the phase spectrum of $U_{\varepsilon}$ 
	(Appendix B, equations \ref{eq:potential-U-epsilon-1-phase-spectrum} and \ref{eq:potential-U-epsilon-2-phase-spectrum}).
	\label{def:eql-min-phase}
\end{definition}
\begin{proposition}
	\label{prop:minimum-phase-change-layers}
	An equivalent layer $p_{\varepsilon}$ having minimum phase changes (Definition \ref{def:eql-min-phase})
	depends only on the spatial distribution of the true sources and not on the particular observed 
	potential field $d^{o}$.
\end{proposition}
\begin{proof}
	It follows from equations \ref{eq:p-phase-spectrum-U} and definition of phase spectrum (Appendix B) that 
	$\mathcal{F} [ p_{\varepsilon}(\mathbf{r}'') ]$ must be the product of a real function and
	$\mathcal{F} [ U_{\varepsilon}(\mathbf{r}^{\dagger}) ]$.
	From equations \ref{eq:potential-U-epsilon-Fourier-layer} and \ref{eq:p-Fourier}, we can verify that 
	$\mathcal{F} [ p_{\varepsilon}(\mathbf{r}'') ]$ is in fact the product of 
	$\mathcal{F} [ U_{\varepsilon}(\mathbf{r}^{\dagger}) ]$ 
	and a term $\left( \nicefrac{\alpha_{d}}{\alpha_{a}} \right) k \, e^{\, k \, (z''_{0} - z^{\dagger}_{0})}$. 
	Since $\nicefrac{\alpha_{d}}{\alpha_{a}}$ is the only term that can be a complex function, we conclude that 
	the imposed condition can only be satisfied for the case in which $\alpha_{a} = \alpha_{r} \, \alpha_{d}$, 
	where $\alpha_{r}$ is a generic real function satisfying $\alpha_{r} \ne 0$ for $0 < k \le k_{\varepsilon}$.
	It can then be verified that, in this case, \ref{eq:p-Fourier} assumes the form
	\begin{equation}
		\mathcal{F} [ p_{\varepsilon}(\mathbf{r}'') ] =
		\tfrac{k}{\alpha_{r}} \, \mathcal{F} [ U_{\varepsilon}(\mathbf{r}^{\dagger}) ] 
		e^{\, k \, (z''_{0} - z^{\dagger}_{0})} \: ,
		\label{eq:p-Fourier-min-phase}
	\end{equation}
	which is the product of a real function and $\mathcal{F} [ U_{\varepsilon}(\mathbf{r}^{\dagger}) ]$.
	Since $\alpha_{r}$ does not depend on $\alpha_{d}$, $\mathcal{F} [ p_{\varepsilon}(\mathbf{r}'') ]$ does not 
	depend on $\alpha_{d}$ as well. This implies that the phase spectrum of $p_{\varepsilon}$ approximates that 
	of potential $U$ (Appendix B), thereby corresponding to an approximation 
	of the phase spectrum associated with the spatial distribution of the (shallow parts of the) true sources.
\end{proof}

Equivalent layers having minimum phase changes (Definition \ref{def:eql-min-phase}) may present intricate 
spatial domain representations, dictated by the particular form of the real function $\alpha_{r}$
(equation \ref{eq:p-Fourier-min-phase}).
Such layers give raise to potential-field transformations (equation \ref{eq:transformed-field-continuous-layer-Fourier})
defined by $\alpha_{t} = \alpha_{r} \, \alpha_{f}$ (equation \ref{eq:alpha-t-generic}), which may also
have non-trivial representations in the space domain.

\subsection{Neumann's equivalent layers}

\begin{definition}
	An equivalent layer $p_{\varepsilon}$ (Definition \ref{def:continuous-eqlayer}) is called Neumann's equivalent 
	layer if it has the following form in the Fourier domain
	\begin{equation}
		\mathcal{F} [ p_{\varepsilon}(\mathbf{r}'') ] =
		k \, \mathcal{F} [ U_{\varepsilon}(\mathbf{r}'') ] \: .
		\label{eq:p-Fourier-Neumann}
	\end{equation}
	\label{def:eql-Neumann}
\end{definition}
\begin{proposition}
	All Neumann's equivalent layers have minimum phase changes.
	\label{prop:eql-Neumann}
\end{proposition}
\begin{proof}
	Equation \ref{eq:p-Fourier-Neumann} is directly obtained from \ref{eq:p-Fourier}, by using 
	$\alpha_{a} = \alpha_{d}$. It can be easily verified that, in this case, 
	$\mathcal{F} [ p_{\varepsilon}(\mathbf{r}'') ]$ is a particular case of equation \ref{eq:p-Fourier-min-phase},
	with $\alpha_{r} = 1$. Consequently, Neumann's equivalent layers have minimum phase changes
	according to Definition \ref{def:eql-min-phase}.
\end{proof}

The predicted potential field $d_{\varepsilon}$ (equation \ref{eq:predicted-data-continuous-layer}) 
due to Neumann's equivalent layers (Definition \ref{def:eql-Neumann}) assumes the form
\begin{equation}
	d_{\varepsilon}(\mathbf{r})
	= \int\limits_{-\infty}^{\infty}\int\limits_{-\infty}^{\infty}
	p_{\varepsilon}(\mathbf{r}'') \: D_{d} \Psi(\mathbf{r}, \mathbf{r}'') \:\: d\mathcal{S}''
	\: , \quad z < z''_{0} \: ,
	\label{eq:predicted-data-Neumann-layer}
\end{equation}
where $D_{d}$ is the same space-domain differential operator associated with
the observed potential field data $d^{o}$ (equation \ref{eq:observed-data-space}).
As a consequence, Neumann's equivalent layers (Definition \ref{def:eql-Neumann}) preserve the physical
relationship between the observed potential field and the true sources.
It means that, if the observed potential field $d^{o}$ is the $x$-component of the 
gravitational attraction due to the true sources, the predicted potential field $d_{\varepsilon}$  
represents the $x$-component of the gravitational attraction due to the planar equivalent layer, for example.
The same is valid for any other potential field.

Any particular Neumann's layer $p_{\varepsilon}$ yields transformed potential 
fields (equation \ref{eq:transformed-field-continuous-layer-Fourier}) 
defined by $\alpha_{t} = \alpha_{f}$ (equation \ref{eq:alpha-t-generic}) in the Fourier domain. 
Hence, the corresponding transformed fields $f_{\varepsilon}$ in the space domain 
(equation \ref{eq:transformed-field-continuous-layer}) assume the simple form
\begin{equation}
	f_{\varepsilon}(\mathbf{r})
	= \int\limits_{-\infty}^{\infty}\int\limits_{-\infty}^{\infty}
	p_{\varepsilon}(\mathbf{r}'') \: D_{f} \Psi(\mathbf{r}, \mathbf{r}'') \:\: d\mathcal{S}''
	\: , \quad z < z''_{0} \: ,
	\label{eq:transformed-field-continuous-layer-Neumann}
\end{equation}
where $D_{f}$ is the space-domain representation of the differential operator $\alpha_{f}$ 
associated with the desired transformed field $f$ (equation \ref{eq:transformed-field-space}).

It follows immediately from equation \ref{eq:p-Fourier-Neumann} that Neumann's equivalent layers approximate 
the vertical derivative $\tfrac{\partial}{\partial z} U(\mathbf{r}'')$ (equation \ref{eq:potential-U-1st-deriv-space}) 
on the layer plane $z''_{0}$.
For the cases in which the observed potential field $d^{o}$ represents gravity data, 
$p_{\varepsilon}$ approximates the gravity disturbance (or vertical component of the gravitational
attraction) due to the true sources on the layer plane $\mathcal{S}''$.
If $d^{o}$ represents magnetic data produced by sources with constant total magnetization direction, 
then $p_{\varepsilon}$ approximates the pseudogravity 
anomaly \citep[][ p. 345]{blakely1996} due to the true sources on $\mathcal{S}''$.
In this case, $p_{\varepsilon}$ is strictly positive because it is defined in terms of a strictly positive 
function $\sigma(\mathbf{r}')$ (equation \ref{eq:potential-U-space}) representing the 
total-magnetization intensity at all points $\mathbf{r}' = (x', y', z')^{\top} \in \mathcal{V}$ 
within the sources.
This result is in perfect agreement with that deduced by \citet{reis_etal2020} for
the particular case in which the observed potential field $d^{o}$ is
the approximated total-field anomaly (equation \ref{eq:potential-U-oblique-2nd-deriv-space}).

It is widely known that integrating the potential field $\tfrac{\partial}{\partial z} U(\mathbf{r}'')$ 
(equation \ref{eq:potential-U-1st-deriv-space}) on a horizontal plane above the causative sources
gives their total mass (or total dipole moment), according to the
Gauss' Theorem \citep[e.g.,][]{hammer1945, lafehr1965, bhattacharyya1967}. 
Furthermore, \citet{medeiros-silva1995a} show that the horizontal first-order moments, second-order moment $xy$, 
difference between the second-order moments $xx$ and $yy$, as well as the horizontal coordinates of the sources center of mass
(or center of dipole moment) of the true sources can also be estimated from the $\tfrac{\partial}{\partial z} U(\mathbf{r}'')$
that they produce on a horizontal plane above them.
Thus, all these parameters of the true sources can be directly estimated from 
Neumann's equivalent layers $p_{\varepsilon}$ (Definition \ref{def:eql-Neumann}) because
they approximate the field $\tfrac{\partial}{\partial z} U(\mathbf{r}'')$ (equation \ref{eq:potential-U-1st-deriv-space}) 
on $\mathcal{S}''$.